% ============================================================
% sections/conclusion.tex  おわりに（初稿 2026-10-26，模擬査読を受けて主張の強さを直した版 2026-10-26 7:30）
% ============================================================
\section{おわりに}\label{sec:conclusion}
本研究は，自他で重みを共有する予測の仕組みである重ね合わせ機構\cite{noguchi2022superposition}が，自分の価値の経験から他者について何を生み，何を生まないかを，価値に従う行動の規則を選ぶ確率と自他の共有の有無を操作した計算機実験で調べた．

与えた快・不快の信号のもとで価値に従って動いた経験は，報酬も目標の入力も用いない予測学習の中で，価値に従う行動の見込みを生み，その強さは経験の量に応じて段階的に強まった．
その見込みは，自分側で学んだ重みを通して他者側に入れた快・不快にも働き，共有だけを外すと他者側の働きだけが消えた．
見込みを用いて研究者が用意した二つの好みの候補を入れ比べると，経験が多いときには偶然を上回って好みの方を選べたが，主の基準（自分の視界の予測誤差）での正解率は，他者の進んだ方角だけで選ぶ基準を上回らず，副の基準（他者側の特徴予測の誤差）でも組をまたいでわずかに上回るにとどまった．
一方で，行動の仕方が自分と違う他者では，本当の快・不快を与えたときの予測が経験の少ない条件より外れ，他者の行動の仕方を自分にそろえると，この悪化は小さくなった．
また，他者用のモジュールは調べた4通りの構成のどれでも他者の好みを出力に書き出さず，予測に役立つ位置や方角の手がかりを伝えた．

これらの結果は，シミュレーションによる他者理解のできることとできないことを，事前に学習した自他の対応を持つ一つのモデルの上で示している．
自分の経験を共有された重みで他者に当てはめると，他者を価値の方へ動く者として予測する枠組みが得られる．
しかし，本研究で調べた他者用モジュールと予測学習の条件では，他者に合わせた調整と，他者固有の好みを出力する働きは確認できなかった．また，候補の照合による判別にも位置や方角による説明が残り，好みを独立に推定できることまでは示されなかった．
ただし，経験の量の効果は学習データの位置の分布の変化と分けられておらず，投影の誤りという解釈にも位置の偏りによる説明が残る．
位置の分布をそろえた対照によってこれを確かめること，および好みの順位でしか他者の行動を予測できない課題で結果がどう変わるかを調べることが，次の課題である．
